Se comparan tres modelos idealizados en \(0<x<L\), con \(L>0\): difusión térmica \(u_t=\kappa u_{xx}\), vibración de una cuerda \(u_{tt}=c^2u_{xx}\) y equilibrio térmico sin fuentes \(U''=0\). El campo está normalizado y es adimensional; \(\kappa>0\) tiene unidades de longitud al cuadrado por tiempo y \(c>0\), de longitud por tiempo. En todos los casos, los valores en ambos extremos se mantienen en cero.

Un registro sintético, cuyo tipo de ensayo no está identificado, proporciona
\[
u(x,0)=\sin(\pi x/L),\qquad u_t(x,0)=0.
\]
Para este registro se exige que el campo y las derivadas que aparecen en cada modelo, además de \(u_t\), se prolonguen continuamente hasta \(t=0\). Determine qué modelos son compatibles con la información completa y explique qué papel desempeña la tasa inicial registrada en cada uno. Justifique la compatibilidad que afirme.

Otro registro, identificado como térmico, indica temperatura inicial normalizada \(u(x,0)=1\) para \(0<x<L\) y extremos a cero para todo \(t>0\). Se proponen dos clases de solución: una exige continuidad en el rectángulo cerrado \([0,L]\times[0,T]\), con \(T>0\); la otra exige una solución clásica para \(t>0\), continua hasta los extremos a cada tiempo positivo, y sólo el límite \(u(x,t)\to1\) para cada \(x\) interior fijo cuando \(t\downarrow0\). ¿Qué clase queda descartada por los datos de las esquinas? ¿Qué permite concluir esto sobre el modelo térmico para tiempos positivos?
